%=====================================================================
% RAPPORT PFE - EVE MOBILITY
% Système de surveillance et de sécurité à distance pour scooter électrique
% BEN HAJ ALI Ayoub & BOUTALEB Ibrahim
% ISIMA Mahdia - 2025/2026
%=====================================================================

\documentclass[12pt, a4paper]{report}

%--- Encodage et langue ---
\usepackage[utf8]{inputenc}
\usepackage[T1]{fontenc}
\usepackage[french]{babel}

%--- Mise en page ---
\usepackage[
  top=2.5cm,
  bottom=2.5cm,
  left=3cm,
  right=2.5cm,
  headheight=15pt
]{geometry}

%--- Typographie ---
\usepackage{times}
\usepackage{microtype}
\usepackage{setspace}
\onehalfspacing

%--- Couleurs ---
\usepackage[dvipsnames]{xcolor}
\definecolor{isimablue}{RGB}{46, 117, 182}
\definecolor{darkblue}{RGB}{31, 78, 121}
\definecolor{lightgray}{RGB}{245, 245, 245}
\definecolor{codetextcolor}{RGB}{50, 50, 50}
\definecolor{conclusionbg}{RGB}{240, 255, 240}
\definecolor{conclusionframe}{RGB}{34, 139, 34}

%--- Graphiques et figures ---
\usepackage{graphicx}
\usepackage{float}
\usepackage{wrapfig}

%--- Tableaux ---
\usepackage{booktabs}
\usepackage{tabularx}
\usepackage{array}
\usepackage{multirow}
\usepackage{longtable}
\usepackage{colortbl}

%--- Listes ---
\usepackage{enumitem}
\setlist[itemize]{label=\textbullet, leftmargin=1.5em, itemsep=2pt}
\setlist[enumerate]{leftmargin=1.5em, itemsep=2pt}

%--- Liens hypertexte ---
\usepackage[
  colorlinks=true,
  linkcolor=isimablue,
  urlcolor=isimablue,
  citecolor=isimablue,
  bookmarks=true,
  bookmarksopen=true
]{hyperref}

%--- En-têtes et pieds de page ---
\usepackage{fancyhdr}
\pagestyle{fancy}
\fancyhf{}
\fancyhead[L]{\small\textcolor{isimablue}{\leftmark}}
\fancyhead[R]{\small\textcolor{isimablue}{EVE Mobility}}
\fancyfoot[C]{\small\textcolor{gray}{\thepage}}
\fancyfoot[L]{\small\textcolor{gray}{BEN HAJ ALI Ayoub \& BOUTALEB Ibrahim}}
\renewcommand{\headrulewidth}{0.4pt}
\renewcommand{\headrule}{\hbox to\headwidth{\color{isimablue}\leaders\hrule height \headrulewidth\hfill}}
\renewcommand{\footrulewidth}{0pt}

%--- Style des chapitres ---
\usepackage{titlesec}

\titleformat{\chapter}[display]
  {\normalfont\huge\bfseries\color{isimablue}}
  {\filleft\Large\color{darkblue}Chapitre \thechapter}
  {1ex}
  {\titlerule[2pt]\vspace{1ex}\filright}
  [\vspace{1ex}\titlerule]

\titleformat{\section}
  {\normalfont\large\bfseries\color{darkblue}}
  {\thesection}{1em}{}
  [\vspace{2pt}{\color{isimablue}\hrule height 0.5pt}\vspace{2pt}]

\titleformat{\subsection}
  {\normalfont\normalsize\bfseries\color{darkblue}}
  {\thesubsection}{1em}{}

\titleformat{\subsubsection}
  {\normalfont\normalsize\itshape\color{darkblue}}
  {\thesubsubsection}{1em}{}

\titlespacing*{\chapter}{0pt}{0pt}{30pt}
\titlespacing*{\section}{0pt}{12pt}{6pt}
\titlespacing*{\subsection}{0pt}{10pt}{4pt}

%--- Table des matières ---
\usepackage{tocloft}
\renewcommand{\cftchapfont}{\bfseries\color{isimablue}}
\renewcommand{\cftchappagefont}{\bfseries\color{isimablue}}
\renewcommand{\cftsecfont}{\color{darkblue}}
\renewcommand{\cftsubsecfont}{\small}
\setlength{\cftbeforechapskip}{8pt}

%--- Boîtes et encadrés ---
\usepackage{tcolorbox}
\tcbuselibrary{skins, breakable}

% Encadré d'introduction de chapitre (bleu clair, italique)
\newtcolorbox{chapterintro}{
  colback=blue!5!white,
  colframe=isimablue!60,
  arc=4pt,
  left=10pt, right=10pt, top=8pt, bottom=8pt,
  fontupper=\itshape
}

% Encadré de conclusion de chapitre (vert clair, normal)
\newtcolorbox{chapterconclusion}{
  colback=conclusionbg,
  colframe=conclusionframe!60,
  arc=4pt,
  left=10pt, right=10pt, top=8pt, bottom=8pt,
  fontupper=\normalfont,
  title={\textbf{\textcolor{conclusionframe}{Conclusion du chapitre}}},
  attach boxed title to top left={yshift=-2mm, xshift=4mm},
  boxed title style={colback=conclusionframe!70, colframe=conclusionframe, arc=3pt}
}

% Encadré info générique
\newtcolorbox{infobox}{
  colback=lightgray,
  colframe=isimablue,
  arc=4pt,
  left=8pt, right=8pt, top=6pt, bottom=6pt,
  fonttitle=\bfseries\color{white},
  title style={fill=isimablue}
}

%--- Code source ---
\usepackage{listings}
\lstset{
  basicstyle=\small\ttfamily\color{codetextcolor},
  backgroundcolor=\color{lightgray},
  frame=single,
  rulecolor=\color{isimablue!50},
  breaklines=true,
  breakatwhitespace=true,
  numbers=left,
  numberstyle=\tiny\color{gray},
  tabsize=2,
  showstringspaces=false,
  captionpos=b,
  xleftmargin=10pt,
  framexleftmargin=10pt
}

%--- Numérotation des figures/tableaux par chapitre ---
\usepackage{chngcntr}
\counterwithin{figure}{chapter}
\counterwithin{table}{chapter}

%--- Glossaire / abréviations ---
\usepackage[printonlyused]{acronym}

%--- Mathématiques ---
\usepackage{amsmath}

%--- Divers ---
\usepackage{parskip}
\setlength{\parindent}{0pt}
\setlength{\parskip}{6pt}

%=====================================================================
% DÉBUT DU DOCUMENT
%=====================================================================
\begin{document}

%=====================================================================
% PAGE DE GARDE
%=====================================================================
\begin{titlepage}
  \thispagestyle{empty}
  \pagenumbering{roman}

  \noindent
  \begin{minipage}[t]{0.48\textwidth}
    \small
    République Tunisienne\\
    Ministère de l'Enseignement Supérieur\\
    et de la Recherche Scientifique
  \end{minipage}
  \hfill
  \begin{minipage}[t]{0.48\textwidth}
    \raggedleft\small
    Université de Monastir\\
    Institut Supérieur d'Informatique\\
    et de Mathématiques de Mahdia
  \end{minipage}

  \vspace{0.4cm}
  {\color{isimablue}\rule{\linewidth}{2pt}}
  \vspace{0.2cm}

  \begin{center}
    \textit{\large Projet de fin d'études}

    \vspace{0.8cm}

    {\fontsize{48}{56}\selectfont\textbf{MÉMOIRE}}

    \vspace{0.5cm}
    \textit{Présenté à}

    \vspace{0.3cm}
    {\Large\textbf{Institut Supérieur d'Informatique et de Mathématiques de Mahdia}}

    \vspace{0.5cm}
    \textit{Effectué à}

    \vspace{0.3cm}
    {\Large\textbf{EVE Mobility}}

    \vspace{0.5cm}
    \textit{En vue de l'obtention du diplôme de}

    \vspace{0.3cm}
    Licence en Informatique

    \vspace{0.6cm}
    \textit{Élaboré par :}

    \vspace{0.3cm}
    {\large\textbf{BEN HAJ ALI Ayoub \quad --- \quad BOUTALEB Ibrahim}}

    \vspace{0.8cm}
    {\color{isimablue}\rule{\linewidth}{2pt}}

    \vspace{0.4cm}
    {\Large\textbf{Système de surveillance et de sécurité à distance\\pour scooter électrique}}

    \vspace{0.4cm}
    {\color{isimablue}\rule{\linewidth}{2pt}}
  \end{center}

  \vspace{0.8cm}

  \begin{minipage}[t]{0.5\textwidth}
    \textbf{Encadrante Pédagogique :}\\[4pt]
    \textbf{Encadrant Professionnel :}
  \end{minipage}
  \begin{minipage}[t]{0.5\textwidth}
    Mme Najoua Achoura\\[4pt]
    M. Lassaad Akrout
  \end{minipage}

  \vfill

  \begin{center}
    Année universitaire 2025--2026
  \end{center}
\end{titlepage}

%=====================================================================
% REMERCIEMENTS
%=====================================================================
\chapter*{Remerciements}
\addcontentsline{toc}{chapter}{Remerciements}
\markboth{Remerciements}{Remerciements}
\thispagestyle{fancy}

Ce projet n'aurait pas vu le jour sans la contribution de nombreuses personnes auxquelles
nous tenons à adresser nos remerciements sincères.

Nous remercions en premier lieu notre encadrante pédagogique, \textbf{Mme Najoua Achoura},
dont les retours précis et les séances de suivi régulières ont orienté nos choix
techniques aux moments clés du projet. Sa capacité à identifier rapidement les points
à améliorer nous a évité plusieurs impasses.

Notre gratitude va également à notre encadrant professionnel, \textbf{M. Lassaad Akrout},
qui nous a accordé sa confiance dès le premier jour du stage et nous a laissé la liberté
de proposer des solutions techniques nouvelles pour EVE Mobility. Sa connaissance terrain
des contraintes liées à la gestion de la flotte de scooters a été déterminante pour
orienter le cahier des charges vers des besoins réels.

Nous remercions toute l'équipe d'EVE Mobility pour leur accueil chaleureux et leur
disponibilité à répondre à nos questions techniques au quotidien.

Enfin, nous remercions nos familles pour leur patience et leur soutien constant, particulièrement
durant les phases intenses de développement et de rédaction.

%=====================================================================
% DÉDICACE
%=====================================================================
\chapter*{Dédicace}
\addcontentsline{toc}{chapter}{Dédicace}
\markboth{Dédicace}{Dédicace}
\thispagestyle{fancy}

\vspace{2cm}

\begin{center}
  \textit{À nos parents,}\\[0.5em]
  \textit{pour leur amour inconditionnel, leurs sacrifices et leurs encouragements}\\
  \textit{constants qui ont été le pilier de notre réussite.}\\[1.5em]

  \textit{À nos frères, sœurs et proches,}\\[0.5em]
  \textit{pour leur présence et leur soutien dans les moments difficiles}\\
  \textit{comme dans les moments de joie.}\\[1.5em]

  \textit{À tous nos amis et collègues,}\\[0.5em]
  \textit{qui ont partagé avec nous les défis et les joies de ce parcours académique.}\\[2em]

  \textbf{\textit{Nous vous dédions ce travail avec toute notre affection et notre reconnaissance.}}
\end{center}

%=====================================================================
% RÉSUMÉ
%=====================================================================
\chapter*{Résumé}
\addcontentsline{toc}{chapter}{Résumé}
\markboth{Résumé}{Résumé}
\thispagestyle{fancy}

Ce travail, mené en stage au sein de la société \textbf{EVE Mobility} (Tunis), répond à
un besoin opérationnel concret~: doter la flotte de scooters électriques de la société
d'un système centralisé de télésurveillance et de détection d'anomalies. Nous avons
conçu et réalisé une architecture IoT trois couches~: un firmware embarqué sur
\textbf{ESP32} (C++/Arduino) qui collecte les données de cinq capteurs --- accéléromètre
\textbf{MPU-6050}, module BMS lithium via Bluetooth, capteurs \textbf{TPMS}, et quatre
détecteurs d'ouverture (siège, capots, compartiment batterie) --- et les publie vers
un broker \textbf{HiveMQ Cloud} via MQTT over TLS~1.3~; un backend \textbf{Supabase}
(PostgreSQL + WebSocket Realtime) qui stocke, sécurise et redistribue les données en
temps réel~; et une application mobile \textbf{React Native / Expo}, déployée sur iOS
et Android, offrant un tableau de bord en temps réel, une gestion des alertes et la
configuration à distance des seuils de détection.

\vspace{0.5cm}
\textbf{Mots-clés :} IoT, ESP32, BMS, MQTT, HiveMQ, React Native, Supabase, temps réel,
scooter électrique, MPU-6050, TPMS, détection de chute, sécurité embarquée.

%=====================================================================
% ABSTRACT
%=====================================================================
\chapter*{Abstract}
\addcontentsline{toc}{chapter}{Abstract}
\markboth{Abstract}{Abstract}
\thispagestyle{fancy}

This internship project, conducted at \textbf{EVE Mobility} (Tunis, Tunisia), addresses
a practical operational need: equipping the company's electric scooter fleet with a
unified remote monitoring and anomaly detection platform. We designed and built a
three-layer IoT architecture. At the embedded level, an \textbf{ESP32} microcontroller
runs a C++/Arduino firmware that reads five sensor types --- an \textbf{MPU-6050}
accelerometer, a lithium \textbf{BMS} module over Bluetooth, \textbf{TPMS} tire sensors,
and four tamper contacts (seat, front cowl, battery compartment, lower cowl) --- and
publishes JSON payloads to a \textbf{HiveMQ Cloud} MQTT broker over TLS~1.3. At the
cloud level, \textbf{Supabase} (PostgreSQL + WebSocket Realtime) stores and
redistributes data with row-level security policies. At the presentation level, a
\textbf{React Native / Expo} mobile application, running on both iOS and Android,
displays live dashboards, manages alerts, and enables remote alarm control.

\vspace{0.5cm}
\textbf{Keywords:} IoT, ESP32, BMS, MQTT, HiveMQ, React Native, Supabase, real-time,
electric scooter, MPU-6050, TPMS, fall detection, embedded security.

%=====================================================================
% LISTE DES ABRÉVIATIONS
%=====================================================================
\chapter*{Liste des abréviations}
\addcontentsline{toc}{chapter}{Liste des abréviations}
\markboth{Liste des abréviations}{Liste des abréviations}
\thispagestyle{fancy}

\begin{acronym}[TPMS]
  \acro{IoT}{Internet of Things (Internet des Objets)}
  \acro{BMS}{Battery Management System (Système de gestion de batterie)}
  \acro{ESP32}{Espressif Systems 32-bit microcontroller}
  \acro{TPMS}{Tire Pressure Monitoring System (Système de surveillance de pression des pneus)}
  \acro{MPU}{Motion Processing Unit}
  \acro{BLE}{Bluetooth Low Energy}
  \acro{SOC}{State of Charge (État de charge)}
  \acro{SOH}{State of Health (État de santé)}
  \acro{API}{Application Programming Interface}
  \acro{UUID}{Universally Unique Identifier}
  \acro{SQL}{Structured Query Language}
  \acro{JSON}{JavaScript Object Notation}
  \acro{JSONB}{JSON Binary (format PostgreSQL)}
  \acro{UI}{User Interface (Interface utilisateur)}
  \acro{UX}{User Experience (Expérience utilisateur)}
  \acro{IDE}{Integrated Development Environment}
  \acro{PFE}{Projet de Fin d'Études}
  \acro{ISIMA}{Institut Supérieur d'Informatique et de Mathématiques de Mahdia}
  \acro{UML}{Unified Modeling Language}
  \acro{RLS}{Row Level Security}
  \acro{ORM}{Object-Relational Mapping}
  \acro{MQTT}{Message Queuing Telemetry Transport}
  \acro{TLS}{Transport Layer Security}
  \acro{JWT}{JSON Web Token}
\end{acronym}

%=====================================================================
% TABLE DES MATIÈRES
%=====================================================================
\clearpage
\tableofcontents
\thispagestyle{fancy}

%=====================================================================
% LISTE DES FIGURES
%=====================================================================
\clearpage
\listoffigures
\addcontentsline{toc}{chapter}{Liste des figures}
\thispagestyle{fancy}

%=====================================================================
% LISTE DES TABLEAUX
%=====================================================================
\clearpage
\listoftables
\addcontentsline{toc}{chapter}{Liste des tableaux}
\thispagestyle{fancy}

%=====================================================================
% INTRODUCTION GÉNÉRALE
%=====================================================================
\clearpage
\pagenumbering{arabic}   % <-- Début de la numérotation arabe (1, 2, 3...)

\chapter*{Introduction Générale}
\addcontentsline{toc}{chapter}{Introduction Générale}
\markboth{Introduction Générale}{Introduction Générale}
\thispagestyle{fancy}

Lorsque nous avons intégré la société \textbf{EVE Mobility} pour notre stage de fin
d'études, nous avons constaté un écart important entre les ambitions de l'entreprise
et ses outils de gestion de flotte. Les opérateurs devaient se déplacer physiquement
jusqu'à chaque scooter pour vérifier l'état des batteries, aucun système ne les
avertissait automatiquement en cas de chute ou d'ouverture non autorisée d'un capot,
et il n'existait aucune interface centralisée pour superviser l'ensemble des véhicules
à distance.

EVE Mobility nous a alors confié la mission de combler ces lacunes en développant un
système de télésurveillance complet. Après analyse des contraintes techniques de la
flotte --- microcontrôleur ESP32 déjà présent sur les véhicules, batteries lithium
communicant en Bluetooth, exigence de faible consommation énergétique --- nous avons
conçu une architecture IoT en trois couches~: une couche embarquée (ESP32 + capteurs),
une couche cloud (Supabase + HiveMQ), et une couche mobile (React Native / Expo).

Le système réalisé permet aux opérateurs d'EVE Mobility de consulter en temps réel
l'état de chaque scooter depuis leur smartphone, de recevoir une alerte instantanée
dès qu'une anomalie est détectée (chute, pression de pneu insuffisante, ouverture d'un
capot), et de contrôler le système d'alarme à distance. Toutes les données sont
persistées dans une base PostgreSQL sécurisée et redistribuées par WebSocket pour
assurer une interface réactive sans rechargement manuel.

Ce rapport est organisé en quatre chapitres~:

\begin{itemize}
  \item Le \textbf{premier chapitre} présente le cadre général du projet~: l'entreprise
        d'accueil, les objectifs, l'analyse de l'existant, la problématique et la solution
        proposée, ainsi que la méthodologie de travail adoptée.

  \item Le \textbf{deuxième chapitre} est dédié au Sprint~0~: l'analyse des besoins
        et la spécification du système, incluant les besoins fonctionnels et non fonctionnels,
        la planification Scrum, la modélisation UML et les choix techniques retenus.

  \item Le \textbf{troisième chapitre} décrit la conception du système~: l'architecture
        globale, la modélisation détaillée des données et les diagrammes UML de conception.

  \item Le \textbf{quatrième chapitre} présente la réalisation concrète du système,
        avec les captures d'écran de l'application et les résultats obtenus.
\end{itemize}

Enfin, une conclusion générale viendra synthétiser les apports de ce travail et ouvrir
des perspectives d'amélioration futures.

%=====================================================================
% CHAPITRE 1 : CADRE GÉNÉRAL DU PROJET
%=====================================================================
\chapter{Cadre général du projet}

\begin{chapterintro}
Ce chapitre retrace les premières semaines de notre stage chez EVE Mobility~: l'analyse
du terrain, l'identification des lacunes existantes dans la gestion de la flotte, et la
construction progressive de notre réponse technique. Nous décrivons comment nous avons
délimité le périmètre du projet, choisi notre architecture, et organisé notre travail
avec l'équipe encadrante tout au long du développement.
\end{chapterintro}

%-------------------------------------------------------------------
\section{Présentation de l'entreprise}
%-------------------------------------------------------------------

\subsection{Généralités}

EVE Mobility est une start-up tunisienne spécialisée dans la mobilité électrique urbaine.
Lors de notre arrivée en stage, l'entreprise disposait déjà d'une flotte de scooters
équipés d'ESP32 et de modules BMS lithium communicant en Bluetooth, mais sans aucun
logiciel de supervision centralisée. La société projetait d'industrialiser ses outils
numériques pour accompagner la croissance de sa flotte et réduire les coûts de
maintenance corrective.

Notre stage s'est ainsi inscrit dans une phase charnière pour l'entreprise~: nous avons
eu l'opportunité de concevoir, de A à Z, le premier système de télémétrie d'EVE Mobility
en collaboration directe avec l'équipe technique.

\subsection{Domaines d'activité}

EVE Mobility intervient principalement sur les axes suivants, que nous avons pu observer
au cours du stage~:

\begin{itemize}
  \item La commercialisation et la maintenance de scooters électriques destinés à des
        opérateurs et des particuliers en milieu urbain.
  \item L'intégration de composants électroniques embarqués (ESP32, BMS, capteurs)
        sur les véhicules de la flotte.
  \item Le développement de solutions numériques (applications mobiles, outils de
        gestion) pour faciliter le suivi opérationnel de la flotte.
  \item La recherche de partenariats avec des fournisseurs de composants IoT pour
        réduire les coûts et améliorer la fiabilité du matériel.
\end{itemize}

%-------------------------------------------------------------------
\section{Objectifs du projet}
%-------------------------------------------------------------------

Après avoir rencontré l'équipe d'EVE Mobility et analysé les contraintes de la flotte,
nous avons défini avec l'encadrant professionnel un objectif central~: remplacer les
vérifications manuelles par un flux de données automatique permettant de surveiller
chaque scooter depuis un téléphone, à n'importe quelle heure, sans déplacement terrain.

Pour atteindre cet objectif, nous avons décliné le travail en plusieurs sous-objectifs
techniques concrets~:

\begin{itemize}
  \item Surveiller en temps réel l'état des batteries~: niveau de charge, tension,
        courant, température et alertes de protection BMS.
  \item Détecter automatiquement les chutes du scooter grâce à un accéléromètre
        MPU-6050, avec un seuil de sensibilité configurable.
  \item Superviser la pression des pneus via des capteurs TPMS, avec génération
        d'alertes en cas de pression insuffisante.
  \item Détecter les tentatives de sabotage ou de manipulation non autorisée des
        parties sensibles du scooter.
  \item Centraliser toutes les données sur une plateforme distante accessible via
        une application mobile multiplateforme (iOS, Android).
  \item Permettre la configuration des seuils d'alerte directement depuis
        l'application mobile.
\end{itemize}

%-------------------------------------------------------------------
\section{Analyse de l'existant}
%-------------------------------------------------------------------

Dès le premier jour du stage, nous avons observé que la gestion technique des scooters
d'EVE Mobility s'appuyait sur des processus manuels non outillés. Pour vérifier le
niveau de charge d'une batterie, un opérateur devait brancher un câble de diagnostic
directement sur le BMS et lire les valeurs sur son téléphone via une application
générique Bluetooth non intégrée. Il n'existait aucune base de données historisant
ces relevés, et les événements critiques (ouverture d'un capot, chute d'un véhicule)
passaient inaperçus faute de capteurs actifs.

Cette situation engendrait plusieurs problèmes concrets que nous avons documentés~:

\begin{itemize}
  \item Les opérateurs ne savaient pas, sans se déplacer, si une batterie était en
        cours de décharge anormale ou proche de la limite critique.
  \item Aucun mécanisme n'existait pour détecter une chute du scooter ou une
        tentative d'ouverture non autorisée d'un compartiment, ce qui exposait la
        flotte à des risques de vol et de vandalisme non détectés.
  \item L'absence de données historiques rendait toute analyse de maintenance
        préventive impossible~: les pannes étaient traitées a posteriori,
        augmentant les coûts de réparation.
  \item La configuration de la sensibilité des capteurs (seuil de détection de
        chute, seuil de pression minimale des pneus) n'était pas réalisable à distance.
\end{itemize}

%-------------------------------------------------------------------
\section{Problématique}
%-------------------------------------------------------------------

La question centrale à laquelle nous devions répondre était~: \textit{comment,
à partir du matériel ESP32 déjà présent sur les scooters EVE Mobility, construire un
pipeline de données complet --- de la lecture du capteur jusqu'à l'alerte sur le
téléphone de l'opérateur --- qui soit fiable, sécurisé, et maintenable par une équipe
non spécialisée en embarqué~?}

Cette question nous a conduits à identifier quatre verrous techniques à lever~:

\begin{itemize}
  \item \textbf{Hétérogénéité des capteurs}~: l'ESP32 devait interroger simultanément
        des périphériques aux protocoles très différents --- I2C pour le MPU-6050,
        Bluetooth GATT pour les modules BMS, RF~315~MHz pour les TPMS, et GPIO
        interruption pour les détecteurs tamper --- sans bloquer aucune tâche.
  \item \textbf{Transmission temps réel à faible consommation}~: les données devaient
        arriver dans l'application mobile en moins de 2~secondes tout en minimisant
        la bande passante et la consommation de l'ESP32.
  \item \textbf{Application multiplateforme sans infrastructure backend dédiée}~: EVE
        Mobility ne disposait pas d'un serveur propre, nous devions choisir un service
        cloud géré qui intègre à la fois la base de données, l'authentification et le
        temps réel.
  \item \textbf{Gestion des accès et sécurité}~: le système devant être déployé en
        production, chaque utilisateur ne devait accéder qu'aux données de ses propres
        scooters, avec une authentification robuste.
\end{itemize}

%-------------------------------------------------------------------
\section{Solution proposée}
%-------------------------------------------------------------------

Face aux quatre verrous identifiés, nous avons construit une réponse architecturale
articulée en trois couches, chacune choisie pour sa capacité à répondre aux contraintes
spécifiques d'EVE Mobility.

\subsection{Couche matérielle (Hardware)}

Nous avons retenu l'ESP32 comme nœud central car il était déjà présent sur les scooters
et dispose nativement du Bluetooth, du Wi-Fi et d'un nombre de GPIO suffisant pour
accueillir tous nos capteurs sans circuit additionnel. Il est connecté aux capteurs
suivants~:

\begin{itemize}
  \item Un module \textbf{BMS} (Battery Management System) communiquant via Bluetooth
        pour la lecture des données de chaque batterie~: niveau de charge (SOC),
        tension, courant, température, état de santé (SOH) et alertes de protection.
  \item Un \textbf{accéléromètre MPU-6050} (interface I2C) pour la détection des
        chutes par comparaison d'une valeur d'accélération brute au seuil configurable
        \texttt{fall\_threshold}.
  \item Des capteurs \textbf{TPMS} pour la surveillance de la pression des pneus
        avant et arrière.
  \item Des capteurs \textbf{anti-sabotage (tamper)} positionnés sur les points
        sensibles du scooter~: capot avant, capot inférieur et compartiment
        des batteries.
\end{itemize}

\subsection{Couche backend (Cloud)}

Nous avons choisi \textbf{Supabase} comme backend géré pour éviter à EVE Mobility
de devoir maintenir un serveur dédié. Supabase regroupe dans une seule plateforme
tout ce dont nous avions besoin~:

\begin{itemize}
  \item Un stockage PostgreSQL structuré pour les données de télémétrie et les
        informations batteries, avec des politiques RLS qui garantissent
        qu'un opérateur ne voit que les scooters qui lui sont assignés.
  \item Une couche d'authentification intégrée (email/mot de passe + gestion des
        sessions JWT) qui nous a évité de coder un système d'auth from scratch.
  \item Un mécanisme de diffusion WebSocket (\textbf{Supabase Realtime}) qui pousse
        automatiquement les nouvelles données vers l'application mobile dès leur
        insertion en base, sans que l'application ait à interroger le serveur
        périodiquement.
\end{itemize}

Pour le transport des données entre l'ESP32 et Supabase, nous avons opté pour
\textbf{HiveMQ Cloud} comme broker MQTT, car l'ESP32 ne peut pas ouvrir directement
une connexion WebSocket Supabase depuis une bibliothèque Arduino. MQTT over TLS (port
8883) constitue une solution légère et fiable pour les microcontrôleurs avec accès
Wi-Fi limité.

\subsection{Couche applicative (Mobile)}

Nous avons développé l'application mobile avec \textbf{React Native} et \textbf{Expo},
ce choix nous permettant d'écrire une seule base de code JavaScript déployable sur
iOS et Android, tout en bénéficiant de l'accès aux API natives (notifications, BLE)
via les modules Expo. L'application offre aux opérateurs d'EVE Mobility~:

\begin{itemize}
  \item Visualisation en temps réel de la liste des scooters avec indicateurs
        d'état (en ligne, hors ligne, en charge).
  \item Dashboard détaillé par scooter affichant les données des batteries, les
        lectures des capteurs et les alertes actives.
  \item Configuration des seuils d'alerte (seuil de chute, seuil de pression
        des pneus) directement depuis l'interface.
  \item Système d'authentification sécurisé avec gestion de session persistante.
\end{itemize}

Le tableau~\ref{tab:technologies} récapitule les principales technologies utilisées
dans le projet.

\begin{table}[H]
  \centering
  \caption{Récapitulatif des technologies utilisées}
  \label{tab:technologies}
  \renewcommand{\arraystretch}{1.3}
  \begin{tabularx}{\textwidth}{>{\bfseries}p{2.5cm} p{3.5cm} X}
    \toprule
    \rowcolor{isimablue!20}
    \textbf{Couche} & \textbf{Technologie} & \textbf{Rôle} \\
    \midrule
    \multirow{5}{*}{Hardware}
               & ESP32              & Microcontrôleur principal \\
               & BMS + Bluetooth    & Surveillance des batteries \\
               & MPU-6050 (I2C)     & Détection des chutes \\
               & TPMS               & Surveillance de la pression des pneus \\
               & Capteurs tamper    & Détection anti-sabotage \\
    \midrule
    \multirow{2}{*}{Backend}
               & Supabase           & Base de données PostgreSQL + Realtime WebSocket \\
               & Supabase Auth      & Authentification et gestion des sessions \\
    \midrule
    \multirow{3}{*}{Mobile}
               & React Native       & Framework mobile multiplateforme \\
               & Expo               & Toolchain de build et déploiement \\
               & React Navigation   & Navigation entre les écrans \\
    \bottomrule
  \end{tabularx}
\end{table}

%-------------------------------------------------------------------
\section{Méthodologie de travail}
%-------------------------------------------------------------------

Compte tenu de la durée limitée du stage et du fait que les besoins d'EVE Mobility
évoluaient au fur et à mesure de nos échanges avec M. Akrout, nous avons opté pour
un découpage en sprints de deux semaines inspiré de \textbf{Scrum}. Cette organisation
nous a permis de soumettre des livrables fonctionnels à chaque fin de sprint, de
recueillir les retours de l'encadrant professionnel, et d'ajuster les priorités sans
remettre en cause l'architecture globale.

En pratique, notre déroulement a été le suivant~:

\begin{enumerate}
  \item \textbf{Analyse des besoins} et rédaction du cahier des charges en
        collaboration avec l'encadrant professionnel chez EVE Mobility.
  \item \textbf{Sélection et intégration des composants matériels}~: choix de
        l'ESP32, des capteurs BMS, MPU-6050, TPMS et des détecteurs anti-sabotage.
  \item \textbf{Développement du firmware ESP32} en C++ (framework Arduino) pour
        la collecte et la transmission des données capteurs.
  \item \textbf{Conception et implémentation de la base de données Supabase}~:
        définition du schéma relationnel, configuration du Realtime et de
        l'authentification.
  \item \textbf{Développement de l'application mobile} React Native / Expo~:
        screens, composants, navigation et intégration Supabase.
  \item \textbf{Tests d'intégration} et validation du système de bout en bout,
        depuis la collecte hardware jusqu'à l'affichage temps réel sur
        l'application.
\end{enumerate}

%-------------------------------------------------------------------
\section*{Conclusion}
\addcontentsline{toc}{section}{Conclusion}
%-------------------------------------------------------------------

\begin{chapterconclusion}
Ce chapitre a posé les bases de notre travail~: nous avons décrit la situation de
départ chez EVE Mobility, les lacunes observées sur le terrain, les quatre verrous
techniques que nous avions à lever, et les choix architecturaux que nous avons faits
pour y répondre. L'architecture trois couches (ESP32 + HiveMQ / Supabase / React
Native) et la découpe en sprints constituent le fil conducteur de tous les chapitres
suivants. Le chapitre~2 détaille maintenant comment nous avons formalisé les besoins
et spécifié le système dans le cadre du Sprint~0.
\end{chapterconclusion}

%=====================================================================
% CHAPITRE 2 : SPRINT 0
%=====================================================================
\chapter{Sprint 0 : Analyse des Besoins et Spécifications}

\begin{chapterintro}
Le Sprint~0 a été notre phase d'investigation et de formalisation~: avant d'écrire
la première ligne de code, nous avons listé exhaustivement ce que le système devait
faire, ce qu'il ne devait pas faire, et dans quelles conditions il devrait fonctionner.
Ce chapitre présente le résultat de ce travail~: acteurs, user stories, backlog
priorisé, planning des sprints, modèles UML et justification du stack technique retenu.
\end{chapterintro}

%-------------------------------------------------------------------
\section{Introduction}
%-------------------------------------------------------------------

Avant de commencer à modéliser ou à coder, nous avons consacré les deux premières
semaines du Sprint~0 à interviewer l'équipe d'EVE Mobility et à observer les processus
de gestion de la flotte sur le terrain. Cette phase d'investigation nous a permis de
passer d'une problématique générale (« surveiller les scooters ») à un ensemble de
cas d'utilisation précis, mesurables et hiérarchisés par valeur métier.

La formalisation présentée dans ce chapitre est le résultat direct de ces échanges
avec M. Lassaad Akrout et de l'analyse des contraintes matérielles des véhicules.

%-------------------------------------------------------------------
\section{Objectifs du Projet}
%-------------------------------------------------------------------

\subsection{Objectif Général}

Suite à l'analyse du terrain, nous avons formulé l'objectif central avec M. Akrout~:
permettre à un opérateur EVE Mobility de connaître, depuis son téléphone, l'état
complet d'un scooter en temps réel --- batteries, capteurs, alertes actives --- et
d'y réagir à distance (contrôle alarme, modification de seuils) sans nécessiter de
compétences techniques particulières. La robustesse (pas de perte de données lors
d'une coupure réseau) et la réactivité (moins de 2~secondes d'affichage) constituaient
les deux critères non négociables définis avec l'encadrant.

\subsection{Objectifs Spécifiques}

Les objectifs spécifiques du projet sont~:

\begin{itemize}
  \item \textbf{Surveillance des batteries}~: détection automatique des deux batteries
        via Bluetooth, lecture du niveau de charge, identification de l'état de
        fonctionnement (charge/décharge/veille) et de l'état de présence
        (connectée/proche/perdue).
  \item \textbf{Surveillance du mouvement}~: détection automatique des chutes et
        des mouvements anormaux ou suspects du scooter.
  \item \textbf{Surveillance des pneus}~: réception et affichage des données de
        pression et température via capteurs TPMS.
  \item \textbf{Gestion de l'alarme}~: indication de l'état de l'alarme
        (activée/désactivée) et possibilité de l'activer ou désactiver à distance.
  \item \textbf{Détection des manipulations}~: détection automatique de l'ouverture
        du siège, du capot avant, du compartiment des batteries et du capot inférieur.
  \item \textbf{Gestion centralisée}~: centralisation de toutes les données sur
        une plateforme distante, consultation en temps réel via application mobile,
        et gestion des alertes.
\end{itemize}

%-------------------------------------------------------------------
\section{Périmètre et Contraintes}
%-------------------------------------------------------------------

\subsection{Fonctionnalités Incluses}

Le projet inclut les fonctionnalités suivantes~:

\begin{itemize}
  \item Surveillance des batteries, des pneus et des mouvements.
  \item Détection des anomalies et tentatives de manipulation
        (ouverture siège, capots, compartiment batterie).
  \item Transmission des données vers une plateforme distante.
  \item Consultation et supervision à distance via application mobile.
  \item Gestion centralisée de toutes les alertes.
\end{itemize}

\subsection{Fonctionnalités Exclues}

Les fonctionnalités suivantes sont explicitement exclues du périmètre du projet~:

\begin{itemize}
  \item Le contrôle de la conduite du scooter.
  \item Toute modification du matériel d'origine du scooter au-delà de l'ajout
        des capteurs.
\end{itemize}

\subsection{Contraintes}

\begin{itemize}
  \item \textbf{Faible consommation énergétique}~: le système embarqué doit minimiser
        sa consommation batterie.
  \item \textbf{Transmission optimisée}~: les données sont transmises de manière
        périodique pour réduire la bande passante.
  \item \textbf{Fonctionnement continu}~: le système doit fonctionner en continu
        pendant l'utilisation du scooter.
  \item \textbf{Disponibilité des composants}~: les composants doivent être
        disponibles localement ou via des plateformes internationales fiables.
\end{itemize}

%-------------------------------------------------------------------
\section{Spécification des Acteurs}
%-------------------------------------------------------------------

Pour modéliser le périmètre fonctionnel du système, nous distinguons les acteurs
suivants~:

\begin{itemize}
  \item \textbf{Propriétaire du Scooter}~: utilisateur final accédant à l'application
        mobile pour monitorer son véhicule en temps réel et recevoir les alertes.
  \item \textbf{Administrateur Système}~: entité disposant du niveau de privilège
        le plus élevé, chargée de la supervision globale de la plateforme (gestion
        des utilisateurs, monitoring, sécurité, configuration).
  \item \textbf{Système de Monitoring Embarqué}~: l'ESP32 et ses capteurs, qui
        collectent les données en temps réel et communiquent avec le cloud.
  \item \textbf{Fournisseur de Services Externes}~: services tiers potentiels
        (notifications push, assurance, services de sécurité).
\end{itemize}

%-------------------------------------------------------------------
\section{Analyse des Besoins}
%-------------------------------------------------------------------

\subsection{Besoins Fonctionnels}

Les besoins fonctionnels définissent les fonctionnalités que le système doit fournir.
Ils sont exprimés sous forme de user stories, extraits du cahier des charges, et
priorisés selon leur importance métier.

\begin{table}[H]
  \centering
  \caption{Besoins fonctionnels du système de surveillance}
  \label{tab:fonctionnels}
  \renewcommand{\arraystretch}{1.3}
  \small
  \begin{tabularx}{\textwidth}{|c|p{3cm}|X|c|}
    \hline
    \rowcolor{isimablue!20}
    \textbf{ID} & \textbf{Fonctionnalité} & \textbf{User Story} & \textbf{Priorité} \\
    \hline
    1  & Authentification & En tant qu'utilisateur, je dois m'authentifier (email/mot de passe) pour accéder à l'application. & 1 \\
    \hline
    2  & Détection batteries (BLE) & En tant que système, je dois détecter automatiquement les deux batteries via Bluetooth et établir la communication. & 1 \\
    \hline
    3  & Lecture niveau batterie & En tant que propriétaire, je peux consulter le niveau de charge de chaque batterie en temps réel. & 1 \\
    \hline
    4  & État batterie & En tant que propriétaire, je peux voir si chaque batterie est en charge, décharge ou veille. & 1 \\
    \hline
    5  & Présence batterie & En tant que propriétaire, je peux déterminer si les batteries sont connectées, proches ou perdues. & 1 \\
    \hline
    6  & Détection des chutes & En tant que système, je dois détecter les chutes du scooter via accéléromètre et générer une alerte. & 1 \\
    \hline
    7  & Détection mouvements suspects & En tant que système, je dois détecter les mouvements anormaux (vol) et envoyer une alerte. & 1 \\
    \hline
    8  & Pression des pneus & En tant que propriétaire, je peux consulter la pression des pneus en temps réel via capteurs TPMS. & 1 \\
    \hline
    9  & Température des pneus & En tant que propriétaire, je peux consulter la température des pneus en temps réel. & 1 \\
    \hline
    10 & État alarme & En tant que propriétaire, je peux voir si le système d'alarme est activé ou désactivé. & 1 \\
    \hline
    11 & Contrôle alarme à distance & En tant que propriétaire, je peux activer ou désactiver l'alarme à distance via l'application. & 1 \\
    \hline
    12 & Détection ouverture siège & En tant que système, je dois détecter l'ouverture du siège et générer une alerte. & 1 \\
    \hline
    13 & Détection capot avant & En tant que système, je dois détecter l'ouverture du capot avant et générer une alerte. & 1 \\
    \hline
    14 & Détection compartiment batterie & En tant que système, je dois détecter l'ouverture du compartiment des batteries et générer une alerte. & 1 \\
    \hline
    15 & Détection capot inférieur & En tant que système, je dois détecter l'ouverture du capot inférieur et générer une alerte. & 1 \\
    \hline
    16 & Dashboard de supervision & En tant qu'administrateur, je peux consulter en temps réel toutes les données centralisées. & 1 \\
    \hline
    17 & Gestion des alertes & En tant que propriétaire, je peux consulter, archiver et gérer toutes mes alertes. & 1 \\
    \hline
    18 & Application mobile & En tant que propriétaire, je peux accéder à une application mobile (iOS/Android) affichant l'état du scooter. & 1 \\
    \hline
  \end{tabularx}
\end{table}

\subsection{Besoins Non-Fonctionnels}

Au-delà des fonctionnalités, le système doit respecter un ensemble de contraintes de
qualité qui ont guidé nos choix techniques et architecturaux.

\begin{table}[H]
  \centering
  \caption{Besoins non-fonctionnels du système de surveillance}
  \label{tab:nonfonctionne}
  \renewcommand{\arraystretch}{1.4}
  \small
  \begin{tabularx}{\textwidth}{|p{2.8cm}|X|}
    \hline
    \rowcolor{isimablue!20}
    \textbf{Attribut} & \textbf{Spécification} \\
    \hline
    \textbf{Performance} & Affichage des données en moins de 2 secondes. Réponse aux requêtes API en moins de 200~ms pour 95\,\% des cas. Alertes critiques traitées en moins de 500~ms. \\
    \hline
    \textbf{Disponibilité} & Disponibilité du backend supérieure à 99,5\,\% en conditions normales. L'application mobile doit fonctionner en mode déconnecté avec synchronisation ultérieure. \\
    \hline
    \textbf{Sécurité} & Communications chiffrées via HTTPS/TLS~1.3. Authentification par tokens JWT. Communication BLE avec chiffrement natif Bluetooth~5.0. \\
    \hline
    \textbf{Intégrité des données} & Validation des entrées côté serveur. Filtrage des valeurs aberrantes des capteurs. Vérification des timestamps et élimination des duplications. \\
    \hline
    \textbf{Maintenabilité} & Code organisé en modules indépendants. Migrations Supabase versionnées. Documentation du code (commentaires JSDoc et Doxygen). \\
    \hline
    \textbf{Scalabilité} & Architecture Supabase avec connection pooling. Support d'au minimum 1000 scooters connectés simultanément. \\
    \hline
    \textbf{Accessibilité mobile} & Application optimisée pour écrans de 4'' à 6''. Compatible iOS~13+ et Android~10+. Consommation batterie client inférieure à 5\,\% par heure. \\
    \hline
    \textbf{Compatibilité} & ESP32 compatible Bluetooth~5.0 (portée $>$~30~mètres). Application React Native compilable sans modifications sur iOS et Android. \\
    \hline
  \end{tabularx}
\end{table}

%-------------------------------------------------------------------
\section{Gestion de Projet avec Scrum}
%-------------------------------------------------------------------

\subsection{Rôles Scrum}

Nous avons adapté Scrum au contexte d'un stage en binôme~: des sprints de deux semaines
avec une réunion de revue à la fin de chaque sprint en présence de Mme Achoura et de
M. Akrout. Les rôles ont été attribués dès le démarrage~:

\begin{table}[H]
  \centering
  \caption{Équipe Scrum du projet}
  \label{tab:scrum}
  \renewcommand{\arraystretch}{1.4}
  \begin{tabularx}{\textwidth}{|p{3cm}|X|p{4cm}|}
    \hline
    \rowcolor{isimablue!20}
    \textbf{Rôle} & \textbf{Responsabilités} & \textbf{Acteur} \\
    \hline
    \textbf{Scrum Master} & Faciliter les sprints, lever les obstacles, assurer le respect du processus Scrum. & Mme Najoua Achoura \\
    \hline
    \textbf{Product Owner} & Définir et prioriser le backlog, valider les livrables, représenter les besoins métier. & M. Lassaad Akrout (EVE Mobility) \\
    \hline
    \textbf{Équipe de développement} & Développement, tests et amélioration continue du produit. & BEN HAJ ALI Ayoub, BOUTALEB Ibrahim \\
    \hline
  \end{tabularx}
\end{table}

\subsection{Élaboration du Backlog Produit}

Nous avons priorisé le backlog selon la technique MoSCoW lors d'un atelier de deux
heures avec M. Akrout. Chaque user story a été classée en~: \textbf{M} (indispensable
avant la soutenance), \textbf{S} (importante mais sans impact bloquant), \textbf{C}
(souhaitable si le temps le permet), ou \textbf{W} (reportée à une version future).
Cette session a mis en évidence que la surveillance des batteries et la détection de
chute étaient les deux priorités absolues pour la société.

\begin{table}[H]
  \centering
  \caption{Product Backlog priorisé}
  \label{tab:backlog}
  \renewcommand{\arraystretch}{1.3}
  \small
  \begin{tabularx}{\textwidth}{|c|p{4cm}|X|c|}
    \hline
    \rowcolor{isimablue!20}
    \textbf{ID} & \textbf{User Story / Tâche} & \textbf{Description} & \textbf{MoSCoW} \\
    \hline
    1  & Authentification utilisateur      & Connexion sécurisée avec gestion de session Supabase Auth. & M \\
    \hline
    2  & Intégration ESP32 + capteurs      & Programmation firmware C/C++, communication BLE. & M \\
    \hline
    3  & Détection batteries via BLE       & Auto-découverte et appairage des deux batteries. & M \\
    \hline
    4  & Lecture niveau batterie           & Lecture tension/courant/température de chaque batterie. & M \\
    \hline
    5  & Détection chutes (accéléromètre)  & Implémentation MPU-6050 et algorithme de détection de chute. & M \\
    \hline
    6  & Détection mouvements suspects     & Détection vol et mouvements anormaux. & M \\
    \hline
    7  & Détecteurs ouvertures             & Intégration détecteurs tamper sur 4 points sensibles. & M \\
    \hline
    8  & Capteurs TPMS                     & Intégration capteurs pression et température des pneus. & M \\
    \hline
    9  & Système d'alarme                  & Lecture état alarme et activation/désactivation à distance. & M \\
    \hline
    10 & Dashboard de supervision          & Interface affichant données centralisées en temps réel. & M \\
    \hline
    11 & Application mobile principale     & App React Native avec écrans principaux. & M \\
    \hline
    12 & Alertes en temps réel             & Système de notifications pour anomalies détectées. & M \\
    \hline
    13 & Gestion et archivage alertes      & Interface listing alertes avec archivage. & S \\
    \hline
    14 & Historique des états              & Consultation historique batterie, mouvements, alertes. & S \\
    \hline
    15 & Export de données                 & Génération de rapports PDF/CSV. & C \\
    \hline
    16 & Localisation GPS (optionnelle)    & Intégration GPS si présente sur le scooter. & W \\
    \hline
    17 & Optimisation consommation         & Réduction consommation ESP32 (deep sleep, polling). & S \\
    \hline
    18 & Tests et déploiement final        & Tests d'intégration, déploiement EAS Build iOS/Android. & M \\
    \hline
  \end{tabularx}
\end{table}

\subsection{Planification des Sprints}

Le projet est structuré en quatre sprints afin de respecter la date de soutenance
prévue en juin 2026.

\begin{table}[H]
  \centering
  \caption{Planification des sprints}
  \label{tab:sprints}
  \renewcommand{\arraystretch}{1.5}
  \small
  \begin{tabularx}{\textwidth}{|c|c|X|p{3.5cm}|}
    \hline
    \rowcolor{isimablue!20}
    \textbf{Release} & \textbf{Sprint} & \textbf{Objectif de Sprint} & \textbf{Période} \\
    \hline
    \multirow{4}{*}{Release~1}
      & Sprint~0 & Analyse des besoins, spécifications, choix techniques, modélisation UML. & 01/04 -- 14/04/2026 \\
    \cline{2-4}
      & Sprint~1 & Authentification utilisateur, intégration Supabase Auth, configuration de base du projet. & 15/04 -- 30/04/2026 \\
    \cline{2-4}
      & Sprint~2 & UI Dashboard principal, intégration données temps réel, monitoring batterie et capteurs. & 01/05 -- 15/05/2026 \\
    \cline{2-4}
      & Sprint~3 & Communication BLE avec ESP32, détection chutes, détecteurs tamper, TPMS. & 16/05 -- 30/05/2026 \\
    \hline
    Release~2
      & Sprint~4 & Système d'alertes, notifications, tests d'intégration, déploiement via EAS Build. & 01/06 -- 15/06/2026 \\
    \hline
  \end{tabularx}
\end{table}

%-------------------------------------------------------------------
\section{Modélisation UML}
%-------------------------------------------------------------------

Nous avons produit quatre types de diagrammes UML pour formaliser le système avant le
développement. Ces modèles ont servi de référence commune entre nous, Mme Achoura et
M. Akrout, et nous ont permis de détecter plusieurs incohérences dans nos hypothèses
initiales avant qu'elles ne se traduisent en bugs.

\subsection{Diagramme de Cas d'Utilisation Global}

Le diagramme de cas d'utilisation offre une vue globale des interactions entre les
acteurs et le système. Les principaux cas d'utilisation par acteur sont~:

\begin{itemize}
  \item \textbf{Propriétaire du Scooter}~: s'authentifier, consulter l'état du scooter,
        monitorer les batteries, recevoir les alertes, contrôler l'alarme à distance,
        consulter l'historique des événements.
  \item \textbf{Administrateur}~: s'authentifier, gérer les utilisateurs, monitorer
        la santé du système, consulter les journaux (logs), générer des rapports.
  \item \textbf{Système Embarqué (ESP32)}~: collecter les données capteurs, détecter
        les anomalies, transmettre les données vers Supabase.
\end{itemize}

\textit{Les diagrammes UML détaillés sont présentés dans le chapitre~3 consacré à la
conception du système.}

\subsection{Diagramme de Classes Principal}

Le tableau~\ref{tab:classes} présente les principales classes du système et leurs
attributs.

\begin{table}[H]
  \centering
  \caption{Classes principales du système de surveillance}
  \label{tab:classes}
  \renewcommand{\arraystretch}{1.4}
  \small
  \begin{tabularx}{\textwidth}{|p{2.8cm}|X|p{3cm}|}
    \hline
    \rowcolor{isimablue!20}
    \textbf{Classe} & \textbf{Attributs principaux} & \textbf{Relations} \\
    \hline
    \texttt{Profile} & id, email, full\_name, role, approved, created\_at & 1 $\rightarrow$ N Scooter \\
    \hline
    \texttt{Scooter} & id, name, model, reference, fall\_threshold, created\_at & 1 $\rightarrow$ N Telemetry, Battery, TpmsSensor \\
    \hline
    \texttt{Telemetry} & id, scooter\_id, charging, temp, alarm, starter, fallen, tamper\_points[3], status, recorded\_at & N $\rightarrow$ 1 Scooter \\
    \hline
    \texttt{Battery} & id, scooter\_id, serial\_number, slot, soc & N $\rightarrow$ 1 Scooter \\
    \hline
    \texttt{TpmsSensor} & id, scooter\_id, sensor\_id, wheel\_position, pressure, temperature, created\_at & N $\rightarrow$ 1 Scooter \\
    \hline
    \texttt{InvitedEmail} & email, full\_name, created\_at & (table de gestion des invitations) \\
    \hline
  \end{tabularx}
\end{table}

%-------------------------------------------------------------------
\section{Environnement et Choix Techniques}
%-------------------------------------------------------------------

\subsection{Environnement Matériel}

\subsubsection{Postes de développement}

Les deux membres de l'équipe ont travaillé sur leurs machines personnelles tout au long
du stage~:

\begin{figure}[H]
  \centering
  \begin{minipage}{0.45\textwidth}
    \centering
    \includegraphics[width=0.85\textwidth]{image_rapport/pc-portable-hp-victus-15-fa0014nk-i5-12500h-12-go-rtx-3050-ti-4g.jpg}
    \caption*{\small HP Victus 15 --- i5-12500H / 20 Go RAM / RTX 3050 4 Go}
  \end{minipage}
  \hfill
  \begin{minipage}{0.45\textwidth}
    \centering
    \includegraphics[width=0.65\textwidth]{image_rapport/pc-portable-lenovo-ideapad-gaming-3-i5-11em-8gb-512ssd-3050ti.jpg}
    \caption*{\small Lenovo IdeaPad Gaming 3 --- i5-11 / 16 Go RAM / RTX 3050 4 Go}
  \end{minipage}
  \caption{Postes de développement utilisés}
  \label{fig:hardware}
\end{figure}

\subsubsection{Matériel Embarqué}

Le scooter de test est équipé des composants suivants~:

\begin{itemize}
  \item \textbf{ESP32}~: microcontrôleur dual-core 240~MHz, Bluetooth~5.0, Wi-Fi intégré.
  \item \textbf{BMS}~: système de gestion de batterie lithium, communication UART/I2C.
  \item \textbf{MPU-6050}~: accéléromètre 6 axes (accélération + gyroscope)
        pour la détection de mouvement et de chute.
  \item \textbf{Capteurs TPMS}~: capteurs de pression des pneus communiquant par RF~433~MHz.
  \item \textbf{Détecteurs tamper}~: contacts normalement fermés (NC) pour la
        détection d'ouverture des capots et compartiments.
  \item \textbf{Module RF 315~MHz}~: kit émetteur/récepteur sans fil sur bande libre 315~MHz.
\end{itemize}

\subsection{Environnement Logiciel}

Voici les outils que nous avons retenus et utilisés quotidiennement~:

\begin{figure}[H]
  \centering
  \renewcommand{\arraystretch}{1.8}
  \begin{tabular}{cccc}
    \includegraphics[height=1.8cm]{image_rapport/de6d51395aee9e6d67ed425ce6bfe683.jpg} &
    \includegraphics[height=1.8cm]{image_rapport/fc97e5a86e33afbc91c41c05d5553a8d.jpg} &
    \includegraphics[height=1.8cm]{image_rapport/f4e29a346c876477075775fa07a1b645.jpg} &
    \includegraphics[height=1.8cm]{image_rapport/e600ac030026406aabce78051a1b082b.jpg} \\
    \small Visual Studio Code & \small GitHub & \small Supabase & \small Arduino IDE \\[0.5em]
    \includegraphics[height=1.8cm]{image_rapport/3447628dee741173e9e91f03ebb2bdc1.jpg} &
    \includegraphics[height=1.8cm]{image_rapport/013fc0afd7e180d65bbce0d60a4fbc28.jpg} & & \\
    \small Expo Go & \small StarUML & & \\
  \end{tabular}
  \caption{Outils logiciels utilisés dans le projet}
  \label{fig:tools}
\end{figure}

\subsection{Technologies et Langages Utilisés}

Le tableau~\ref{tab:stack} récapitule le stack technique que nous avons retenu, en
précisant pour chaque technologie la raison concrète qui a motivé notre choix dans
le contexte du projet EVE Mobility.

\begin{table}[H]
  \centering
  \caption{Stack technique du projet}
  \label{tab:stack}
  \renewcommand{\arraystretch}{1.4}
  \small
  \begin{tabularx}{\textwidth}{|p{2.5cm}|p{3.5cm}|X|}
    \hline
    \rowcolor{isimablue!20}
    \textbf{Couche} & \textbf{Technologie} & \textbf{Justification} \\
    \hline
    \multirow{4}{*}{Mobile}
      & React Native (Expo)  & Framework multiplateforme iOS/Android depuis une base de code unique. \\
    \cline{2-3}
      & React Navigation     & Gestion de la navigation (Stack, Tab Navigator). \\
    \cline{2-3}
      & JavaScript           & Langage principal de l'application mobile. \\
    \cline{2-3}
      & React Hooks          & Gestion d'état moderne (useState, useEffect, useContext). \\
    \hline
    \multirow{3}{*}{Backend}
      & Supabase             & BaaS basé sur PostgreSQL avec Auth, Realtime WebSocket. \\
    \cline{2-3}
      & PostgreSQL           & Base de données relationnelle fiable avec RLS. \\
    \cline{2-3}
      & HiveMQ Cloud         & Broker MQTT cloud sécurisé (TLS 1.3, port 8883). \\
    \hline
    \multirow{3}{*}{Embarqué}
      & C/C++ Arduino        & Programmation native ESP32 pour performances optimales. \\
    \cline{2-3}
      & PubSubClient         & Bibliothèque MQTT pour l'ESP32. \\
    \cline{2-3}
      & I2C / Wire           & Protocole filaire vers MPU-6050. \\
    \hline
    \multirow{3}{*}{Protocoles}
      & MQTT over TLS        & Transmission sécurisée des données IoT vers le cloud. \\
    \cline{2-3}
      & HTTPS / TLS~1.3      & Chiffrement des communications client-serveur. \\
    \cline{2-3}
      & WebSocket            & Communication bidirectionnelle temps réel via Supabase Realtime. \\
    \hline
  \end{tabularx}
\end{table}

%-------------------------------------------------------------------
\section*{Conclusion}
\addcontentsline{toc}{section}{Conclusion}
%-------------------------------------------------------------------

\begin{chapterconclusion}
Au terme du Sprint~0, nous disposions d'un cahier des charges complet validé avec EVE
Mobility~: 18 user stories priorisées, un backlog de 18 items MoSCoW, un planning sur
5 sprints, et un stack technique justifié poste par poste. Cette phase d'investigation
de deux semaines s'est révélée indispensable~: elle nous a permis de démarrer le
développement sans ambiguïté sur les fonctionnalités attendues, et d'éviter plusieurs
changements de direction coûteux qui auraient compromis les délais de livraison.
\end{chapterconclusion}

%=====================================================================
% CHAPITRE 3 : CONCEPTION DU SYSTÈME
%=====================================================================
\chapter{Conception du Système}

\begin{chapterintro}
Ce chapitre décrit les décisions de conception que nous avons prises avant d'écrire
le code~: comment les trois couches communiquent entre elles, comment nous avons
structuré les données dans PostgreSQL, et comment chaque flux critique (détection de
chute, authentification, commande d'alarme) se déroule étape par étape. La plupart
de ces décisions ont émergé de contraintes très concrètes découvertes lors des
premières expérimentations sur le matériel réel.
\end{chapterintro}

%-------------------------------------------------------------------
\section{Architecture Globale du Système}
%-------------------------------------------------------------------

L'architecture que nous avons conçue répond à une contrainte forte~: l'ESP32 ne peut
pas ouvrir directement une connexion WebSocket vers Supabase depuis les bibliothèques
Arduino. Nous avons donc introduit un broker MQTT (HiveMQ Cloud) comme couche
intermédiaire, ce qui découple l'embarqué du cloud et simplifie la reconnexion
automatique après une coupure réseau. L'architecture résultante comporte trois couches~:

\begin{itemize}
  \item \textbf{Couche embarquée (Edge)}~: l'ESP32 lit les cinq types de capteurs
        et publie des messages JSON compacts vers HiveMQ Cloud via MQTT over TLS~1.3
        (port 8883). Une fonction Supabase Edge Function s'abonne au broker et insère
        les données reçues dans les tables PostgreSQL.
  \item \textbf{Couche cloud (Backend)}~: Supabase stocke les données dans des tables
        relationnelles sécurisées par RLS, authentifie les utilisateurs, et diffuse
        chaque insertion en temps réel vers les clients abonnés via WebSocket Realtime.
  \item \textbf{Couche présentation (Mobile)}~: l'application React Native s'abonne
        aux canaux Supabase Realtime de chaque scooter, met à jour l'UI à chaque
        nouvelle donnée reçue, et émet des commandes (alarme) via des mises à jour
        REST vers Supabase qui les redirigent vers l'ESP32 via MQTT.
\end{itemize}

\begin{figure}[H]
  \centering
  \includegraphics[width=0.9\textwidth]{image_rapport/architecture_globale.png}
  \caption{Architecture globale du système de surveillance}
  \label{fig:architecture}
\end{figure}

%-------------------------------------------------------------------
\section{Modélisation UML}
%-------------------------------------------------------------------

\subsection{Diagramme de Cas d'Utilisation Détaillé}

\subsubsection{Cas d'utilisation — Propriétaire du Scooter}

Le propriétaire est l'acteur principal de l'application mobile. Ses cas d'utilisation
sont~:

\begin{itemize}
  \item \textbf{S'authentifier}~: connexion par email/mot de passe via Supabase Auth,
        avec gestion de session persistante (JWT).
  \item \textbf{Consulter le tableau de bord}~: visualisation en temps réel de
        l'état global du scooter.
  \item \textbf{Surveiller les batteries}~: consultation SOC, tension, courant,
        température et état de fonctionnement.
  \item \textbf{Surveiller les pneus}~: lecture pression et température via TPMS.
  \item \textbf{Consulter les alertes}~: accès à la liste des alertes actives et
        archivées.
  \item \textbf{Contrôler l'alarme}~: activation/désactivation à distance.
  \item \textbf{Configurer les seuils}~: paramétrage du seuil de détection de chute
        et des seuils de pression.
\end{itemize}

\subsubsection{Cas d'utilisation — Administrateur Système}

\begin{itemize}
  \item \textbf{Gérer les utilisateurs}~: création, modification et suppression
        des comptes (accès sur invitation uniquement).
  \item \textbf{Superviser la flotte}~: monitoring centralisé de l'ensemble des
        scooters enregistrés.
  \item \textbf{Consulter les journaux}~: accès aux événements système et aux
        historiques de connexion.
  \item \textbf{Générer des rapports}~: export des données de télémétrie.
\end{itemize}

\subsubsection{Cas d'utilisation — Système Embarqué (ESP32)}

\begin{itemize}
  \item \textbf{Collecter les données}~: lecture périodique MPU-6050 (500~ms),
        BMS (BLE), TPMS (RF), détecteurs tamper (interruption).
  \item \textbf{Détecter les anomalies}~: chute ($M = \sqrt{x^2 + y^2 + z^2}$),
        mouvement suspect, ouverture des capots.
  \item \textbf{Transmettre les données}~: publication JSON vers Supabase via MQTT.
  \item \textbf{Recevoir les commandes}~: écoute du topic MQTT pour les commandes
        d'alarme.
\end{itemize}

\begin{figure}[H]
  \centering
  \includegraphics[width=0.85\textwidth]{image_rapport/use_case_global.png}
  \caption{Diagramme de cas d'utilisation global}
  \label{fig:usecase}
\end{figure}

%-------------------------------------------------------------------
\subsection{Diagramme de Classes}
%-------------------------------------------------------------------

\begin{figure}[H]
  \centering
  \includegraphics[width=0.9\textwidth]{image_rapport/class_diagram.png}
  \caption{Diagramme de classes du système}
  \label{fig:classes}
\end{figure}

\begin{table}[H]
  \centering
  \caption{Description détaillée des classes du système}
  \label{tab:classes_detail}
  \renewcommand{\arraystretch}{1.4}
  \small
  \begin{tabularx}{\textwidth}{|p{2.5cm}|X|p{3.5cm}|}
    \hline
    \rowcolor{isimablue!20}
    \textbf{Classe} & \textbf{Attributs et méthodes} & \textbf{Relations} \\
    \hline
    \texttt{Profile} &
      \textit{Attributs~:} id (UUID), email, full\_name, role (admin/user),
      approved (bool), created\_at\newline
      \textit{Méthodes~:} isAdmin(), isApproved() &
      (lié à \texttt{auth.users} Supabase) \\
    \hline
    \texttt{Scooter} &
      \textit{Attributs~:} id (UUID), name, model, reference,
      fall\_threshold (float), created\_at\newline
      \textit{Méthodes~:} getLatestTelemetry(), getBatteries(), getSensors() &
      1 $\rightarrow$ N \texttt{Telemetry}\newline
      1 $\rightarrow$ N \texttt{Battery}\newline
      1 $\rightarrow$ N \texttt{TpmsSensor} \\
    \hline
    \texttt{Telemetry} &
      \textit{Attributs~:} id, scooter\_id, charging (bool), temp (°C),
      alarm (bool), starter (bool), fallen (bool),
      tamper\_points (bool[3]), status (text), recorded\_at\newline
      \textit{Méthodes~:} hasFallen(), hasTamper() &
      N $\rightarrow$ 1 \texttt{Scooter} \\
    \hline
    \texttt{Battery} &
      \textit{Attributs~:} id, scooter\_id, serial\_number, slot (1 ou 2),
      soc (\%)\newline
      \textit{Méthodes~:} getChargeLevel(), isLow() &
      N $\rightarrow$ 1 \texttt{Scooter} \\
    \hline
    \texttt{TpmsSensor} &
      \textit{Attributs~:} id, scooter\_id, sensor\_id, wheel\_position
      (AV/AR), pressure (bar), temperature (°C), created\_at\newline
      \textit{Méthodes~:} checkPressure() &
      N $\rightarrow$ 1 \texttt{Scooter} \\
    \hline
    \texttt{InvitedEmail} &
      \textit{Attributs~:} email (PK), full\_name, created\_at\newline
      \textit{Méthodes~:} invite(), revoke() &
      (table d'invitation admin) \\
    \hline
  \end{tabularx}
\end{table}

%-------------------------------------------------------------------
\subsection{Diagrammes de Séquence}
%-------------------------------------------------------------------

\subsubsection{Séquence : Détection et Remontée d'une Chute}

\begin{enumerate}
  \item L'ESP32 lit une seule valeur d'accélération brute depuis le MPU-6050 via
        I2C (registre 0x3B) toutes les 500~ms. Pas de calcul complexe : une seule
        valeur en $g$.
  \item Il compare la valeur absolue au seuil \texttt{fall\_threshold} stocké
        dans la table \texttt{scooters} (par défaut~: 2,5~g). Si $|\text{accel}| >
        \text{fall\_threshold}$, alors \texttt{fallen = true}.
  \item L'ESP32 publie un message JSON sur le topic MQTT avec les champs
        \texttt{"accel": 3.2, "fallen": true}.
  \item Une Edge Function Supabase reçoit le message et insère une ligne dans la
        table \texttt{telemetry} avec les valeurs parsées.
  \item Supabase Realtime pousse automatiquement cette insertion vers tous les
        clients abonnés via WebSocket.
  \item L'application React Native reçoit l'événement et affiche une alerte
        critique si \texttt{fallen = true}.
\end{enumerate}

\begin{figure}[H]
  \centering
  \includegraphics[width=0.9\textwidth]{image_rapport/seq_fall.png}
  \caption{Diagramme de séquence~: détection et remontée d'une chute}
  \label{fig:seq_fall}
\end{figure}

\subsubsection{Séquence : Authentification Utilisateur}

\begin{enumerate}
  \item L'utilisateur saisit ses identifiants dans l'écran de connexion.
  \item L'application appelle \texttt{supabase.auth.signInWithPassword()}.
  \item Supabase Auth vérifie les credentials et retourne un token JWT.
  \item L'application stocke le token et redirige vers le dashboard.
  \item À chaque requête suivante, le token JWT est joint automatiquement dans
        l'en-tête \texttt{Authorization: Bearer <token>}.
\end{enumerate}

\begin{figure}[H]
  \centering
  \includegraphics[width=0.85\textwidth]{image_rapport/seq_auth.png}
  \caption{Diagramme de séquence~: authentification utilisateur}
  \label{fig:seq_auth}
\end{figure}

\subsubsection{Séquence : Contrôle de l'Alarme à Distance}

\begin{enumerate}
  \item L'utilisateur appuie sur le bouton d'activation/désactivation de l'alarme
        dans l'application mobile.
  \item L'application insère une nouvelle ligne dans \texttt{telemetry} avec le
        champ \texttt{alarm} positionné à \texttt{true} ou \texttt{false}.
  \item Supabase Realtime publie ce changement vers tous les abonnés au canal
        du scooter concerné.
  \item L'ESP32, abonné au topic MQTT \texttt{scooter/<id>/commands} via HiveMQ,
        reçoit la commande et active ou coupe physiquement le relais de l'alarme.
  \item L'ESP32 confirme en publiant l'état effectif (\texttt{alarm}) dans son
        prochain message de télémétrie sur \texttt{scooter/<id>/telemetry}.
  \item L'interface de l'application se met à jour en temps réel via Realtime.
\end{enumerate}

\begin{figure}[H]
  \centering
  \includegraphics[width=0.85\textwidth]{image_rapport/seq_alarm.png}
  \caption{Diagramme de séquence~: contrôle de l'alarme à distance}
  \label{fig:seq_alarm}
\end{figure}

%-------------------------------------------------------------------
\subsection{Diagramme d'États du Scooter}
%-------------------------------------------------------------------

Les états du scooter sont dérivés des colonnes de la table \texttt{telemetry}~:

\begin{itemize}
  \item \textbf{Hors ligne (offline)}~: champ \texttt{status = 'offline'} --- aucune
        donnée reçue depuis plus de 60~secondes. L'application affiche le scooter
        en gris.
  \item \textbf{En ligne (online)}~: champ \texttt{status = 'online'} --- télémétrie
        reçue régulièrement, aucune anomalie (\texttt{fallen = false},
        \texttt{alarm = false}, \texttt{tamper\_points = \{f,f,f\}}).
  \item \textbf{En charge (Charging)}~: \texttt{charging = true} --- signalé par
        le firmware lorsque le BMS détecte un courant de charge entrant.
  \item \textbf{Chute détectée (Fallen)}~: \texttt{fallen = true} --- l'accélération
        brute lue du MPU-6050 a dépassé le seuil \texttt{fall\_threshold},
        déclenche une notification critique dans l'application.
  \item \textbf{Alarme active (Alarm)}~: \texttt{alarm = true} --- le relais
        physique d'alarme est activé sur le scooter.
  \item \textbf{Sabotage (Tamper)}~: au moins un élément de \texttt{tamper\_points}
        est \texttt{true} --- ouverture d'un des trois capots surveillés.
\end{itemize}

\begin{figure}[H]
  \centering
  \includegraphics[width=0.85\textwidth]{image_rapport/state_diagram.png}
  \caption{Diagramme d'états du scooter}
  \label{fig:states}
\end{figure}

%-------------------------------------------------------------------
\section{Conception de la Base de Données}
%-------------------------------------------------------------------

\subsection{Schéma Relationnel}

\begin{table}[H]
  \centering
  \caption{Schéma de la table \texttt{scooters}}
  \label{tab:table_scooters}
  \renewcommand{\arraystretch}{1.4}
  \small
  \begin{tabularx}{\textwidth}{|p{3cm}|p{2.5cm}|X|}
    \hline
    \rowcolor{isimablue!20}
    \textbf{Colonne} & \textbf{Type} & \textbf{Description} \\
    \hline
    id             & UUID (PK)          & Identifiant unique du scooter (généré automatiquement) \\
    \hline
    name           & TEXT NOT NULL      & Nom affiché du scooter dans l'application \\
    \hline
    model          & TEXT               & Modèle du scooter \\
    \hline
    reference      & TEXT               & Référence interne EVE Mobility \\
    \hline
    fall\_threshold & DOUBLE PRECISION   & Seuil de détection de chute configurable (défaut~: 2,5) \\
    \hline
    created\_at    & TIMESTAMP          & Date d'enregistrement du scooter \\
    \hline
  \end{tabularx}
\end{table}

\begin{table}[H]
  \centering
  \caption{Schéma de la table \texttt{profiles}}
  \label{tab:table_profiles}
  \renewcommand{\arraystretch}{1.4}
  \small
  \begin{tabularx}{\textwidth}{|p{3cm}|p{2.5cm}|X|}
    \hline
    \rowcolor{isimablue!20}
    \textbf{Colonne} & \textbf{Type} & \textbf{Description} \\
    \hline
    id       & UUID (PK, FK)       & Référence vers \texttt{auth.users} Supabase \\
    \hline
    email    & TEXT NOT NULL       & Adresse email de l'utilisateur \\
    \hline
    full\_name & TEXT              & Nom complet de l'utilisateur \\
    \hline
    role     & TEXT NOT NULL       & Rôle~: \texttt{admin} ou \texttt{user} \\
    \hline
    approved & BOOLEAN NOT NULL    & Accès accordé par l'admin (défaut~: false) \\
    \hline
    created\_at & TIMESTAMPTZ      & Date de création du profil \\
    \hline
  \end{tabularx}
\end{table}

\begin{table}[H]
  \centering
  \caption{Schéma de la table \texttt{telemetry}}
  \label{tab:table_telemetry}
  \renewcommand{\arraystretch}{1.4}
  \small
  \begin{tabularx}{\textwidth}{|p{3cm}|p{2.5cm}|X|}
    \hline
    \rowcolor{isimablue!20}
    \textbf{Colonne} & \textbf{Type} & \textbf{Description} \\
    \hline
    id            & UUID (PK)        & Identifiant unique de la mesure \\
    \hline
    scooter\_id   & UUID (FK)        & Référence vers le scooter \\
    \hline
    charging      & BOOLEAN          & Batterie en cours de charge \\
    \hline
    temp          & DOUBLE           & Température ambiante (°C) \\
    \hline
    alarm         & BOOLEAN          & État de l'alarme (activée/désactivée) \\
    \hline
    starter       & BOOLEAN          & Démarreur actif \\
    \hline
    fallen        & BOOLEAN          & Chute détectée par le MPU-6050 \\
    \hline
    tamper\_points & BOOLEAN[3]      & Points de sabotage~: [siège, capot avant, capot inférieur] \\
    \hline
    status        & TEXT             & État global~: \texttt{online} / \texttt{offline} \\
    \hline
    recorded\_at  & TIMESTAMP        & Horodatage de la mesure \\
    \hline
  \end{tabularx}
\end{table}

\begin{table}[H]
  \centering
  \caption{Schéma de la table \texttt{batteries}}
  \label{tab:table_batteries}
  \renewcommand{\arraystretch}{1.4}
  \small
  \begin{tabularx}{\textwidth}{|p{3cm}|p{2.5cm}|X|}
    \hline
    \rowcolor{isimablue!20}
    \textbf{Colonne} & \textbf{Type} & \textbf{Description} \\
    \hline
    id            & UUID (PK)  & Identifiant unique \\
    \hline
    scooter\_id   & UUID (FK)  & Référence vers le scooter \\
    \hline
    serial\_number & TEXT NOT NULL & Numéro de série BMS de la batterie \\
    \hline
    slot          & INTEGER    & Emplacement de la batterie (1 ou 2) \\
    \hline
    soc           & REAL       & État de charge en \% (0--100) \\
    \hline
  \end{tabularx}
\end{table}

\begin{table}[H]
  \centering
  \caption{Schéma de la table \texttt{tpms\_sensors}}
  \label{tab:table_tpms}
  \renewcommand{\arraystretch}{1.4}
  \small
  \begin{tabularx}{\textwidth}{|p{3cm}|p{2.5cm}|X|}
    \hline
    \rowcolor{isimablue!20}
    \textbf{Colonne} & \textbf{Type} & \textbf{Description} \\
    \hline
    id             & UUID (PK)   & Identifiant unique \\
    \hline
    scooter\_id    & UUID (FK)   & Référence vers le scooter \\
    \hline
    sensor\_id     & TEXT NOT NULL & Identifiant matériel du capteur TPMS \\
    \hline
    wheel\_position & TEXT       & Position~: \texttt{AV} (avant) ou \texttt{AR} (arrière) \\
    \hline
    pressure       & REAL        & Pression en bar \\
    \hline
    temperature    & REAL        & Température en °C \\
    \hline
    created\_at    & TIMESTAMPTZ & Horodatage de la dernière mesure \\
    \hline
  \end{tabularx}
\end{table}

\subsection{Politique de Sécurité (Row Level Security)}

\begin{lstlisting}[language=SQL, caption={Politiques RLS sur les tables principales}]
-- Seul l'utilisateur authentifie voit son propre profil
CREATE POLICY "profiles_self_policy" ON profiles
  FOR ALL USING (auth.uid() = id);

-- Les admins voient tous les profils (gestion des comptes)
CREATE POLICY "profiles_admin_policy" ON profiles
  FOR ALL USING (
    EXISTS (
      SELECT 1 FROM profiles
      WHERE id = auth.uid() AND role = 'admin'
    )
  );

-- Acces aux scooters : uniquement les utilisateurs approuves
CREATE POLICY "scooters_approved_policy" ON scooters
  FOR SELECT USING (
    EXISTS (
      SELECT 1 FROM profiles
      WHERE id = auth.uid() AND approved = true
    )
  );

-- Acces a la telemetrie : meme condition
CREATE POLICY "telemetry_approved_policy" ON telemetry
  FOR ALL USING (
    EXISTS (
      SELECT 1 FROM profiles
      WHERE id = auth.uid() AND approved = true
    )
  );
\end{lstlisting}

%-------------------------------------------------------------------
\section{Conception du Firmware ESP32}
%-------------------------------------------------------------------

\subsection{Architecture du Firmware}

Le firmware de l'ESP32 est structuré autour d'une boucle principale (\texttt{loop()})
qui orchestre plusieurs tâches~:

\begin{itemize}
  \item \textbf{Tâche MQTT}~: maintien de la connexion au broker HiveMQ Cloud et
        traitement des messages entrants.
  \item \textbf{Tâche MPU-6050}~: lecture périodique des angles (500~ms) et
        détection de chute.
  \item \textbf{Tâche Contact/Tamper}~: surveillance des capteurs à contact
        (détection immédiate sur changement d'état).
  \item \textbf{Tâche BLE}~: communication avec les modules BMS.
  \item \textbf{Tâche TPMS}~: décodage des trames RF~315~MHz.
\end{itemize}

\subsection{Format des Messages MQTT}

\begin{lstlisting}[language=,caption={Format JSON — Données MPU-6050 (accélération brute)}]
{
  "type"   : "gyro",
  "accel"  : 2.34,
  "fallen" : false,
  "ts"     : 183420
}
\end{lstlisting}

\begin{lstlisting}[language=,caption={Format JSON — Données capteurs de contact (tamper\_points)}]
{
  "type"          : "contact",
  "tamper_points" : [true, false, false],
  "ts"            : 183470
}
\end{lstlisting}

\begin{lstlisting}[language=,caption={Format JSON — Données de télémétrie générale}]
{
  "type"     : "telemetry",
  "charging" : false,
  "alarm"    : false,
  "starter"  : false,
  "temp"     : 27.4,
  "status"   : "online",
  "ts"       : 183500
}
\end{lstlisting}

%-------------------------------------------------------------------
\section*{Conclusion}
\addcontentsline{toc}{section}{Conclusion}
%-------------------------------------------------------------------

\begin{chapterconclusion}
Les modèles de conception présentés dans ce chapitre ont guidé l'intégralité de notre
développement. La décision d'introduire HiveMQ comme broker intermédiaire --- prise
lors d'une expérimentation qui a montré que l'ESP32 ne pouvait pas maintenir une
connexion WebSocket stable sous Arduino --- a été le pivot de toute l'architecture.
Les politiques RLS définies dès la conception ont permis d'assurer la séparation des
données entre opérateurs sans surcharge applicative. Le chapitre suivant présente la
réalisation concrète et les résultats des tests d'intégration.
\end{chapterconclusion}

%=====================================================================
% CHAPITRE 4 : RÉALISATION DU SYSTÈME
%=====================================================================
\chapter{Réalisation du Système}

\begin{chapterintro}
Ce chapitre raconte le passage du modèle au code~: comment nous avons concrètement
implémenté le firmware ESP32, câblé les capteurs, configuré Supabase et HiveMQ, et
développé chaque écran de l'application mobile. Nous présentons également les résultats
mesurés lors des tests d'intégration que nous avons conduits sur le matériel réel d'EVE
Mobility, ainsi que les difficultés rencontrées et les ajustements apportés en cours
de sprint.
\end{chapterintro}

%-------------------------------------------------------------------
\section{Introduction}
%-------------------------------------------------------------------

Nous avons développé le système en parallèle sur deux fronts~: Ibrahim a pris en charge
le firmware ESP32 et le câblage des capteurs, tandis qu'Ayoub développait les écrans
de l'application mobile et configurait Supabase. Cette organisation nous a permis de
terminer les Sprints~1 et~2 dans les délais. À partir du Sprint~3, les deux parties
ont été intégrées pour les premiers tests bout-en-bout sur le scooter de test mis à
notre disposition par EVE Mobility.

%-------------------------------------------------------------------
\section{Réalisation du Firmware ESP32}
%-------------------------------------------------------------------

\subsection{Environnement de Développement Embarqué}

Nous avons choisi le framework Arduino pour l'ESP32 pour sa maturité et la richesse
de son écosystème de bibliothèques, préféré à ESP-IDF (qui aurait offert plus de
contrôle mais demandé beaucoup plus de temps de développement). Les bibliothèques
tierces intégrées~:

\begin{itemize}
  \item \textbf{PubSubClient}~: client MQTT pour la communication avec HiveMQ Cloud.
  \item \textbf{WiFiClientSecure}~: connexion TLS sécurisée (port 8883).
  \item \textbf{ArduinoJson}~: sérialisation et désérialisation des messages JSON.
  \item \textbf{Wire}~: communication I2C avec le capteur MPU-6050.
\end{itemize}

\subsection{Implémentation du Module MPU-6050}

Pour la détection de chute, nous lisons une seule valeur d'accélération brute du
MPU-6050 via I2C (registre 0x3B), sans calcul complexe de magnitude. Cette approche
réduit drastiquement la charge CPU et la consommation mémoire comparée aux bibliothèques
complètes. Nous comparons simplement cette valeur au seuil configurable~:

\begin{lstlisting}[language=C++, caption={Lecture de l'accélération brute depuis le MPU-6050}]
float readMPUAccel() {
  Wire.beginTransmission(MPU_ADDR);
  Wire.write(0x3B);           // Registre ACCEL_XOUT_H
  Wire.endTransmission(false);
  Wire.requestFrom((uint8_t)MPU_ADDR, (size_t)2, true);

  int16_t accelRaw = Wire.read() << 8 | Wire.read();
  // Conversion : 16384 LSB/g (sensibilité par défaut)
  float accelG = accelRaw / 16384.0;

  return accelG;  // Retourne une seule valeur en g
}
\end{lstlisting}

\subsection{Boucle Principale du Firmware}

\begin{lstlisting}[language=C++, caption={Boucle principale --- Gestion MQTT et capteurs}]
void loop() {
  if (!client.connected()) reconnect();
  client.loop();

  // Capteurs tamper : detection sur changement d'etat
  bool tamper[3] = {
    digitalRead(PIN_SEAT)  == LOW,
    digitalRead(PIN_FRONT) == LOW,
    digitalRead(PIN_LOWER) == LOW
  };
  if (memcmp(tamper, lastTamper, 3) != 0) {
    memcpy(lastTamper, tamper, 3);
    JsonDocument doc;
    doc["type"] = "contact";
    doc["tamper_points"][0] = tamper[0];
    doc["tamper_points"][1] = tamper[1];
    doc["tamper_points"][2] = tamper[2];
    doc["ts"] = millis();
    String payload;
    serializeJson(doc, payload);
    client.publish(mqtt_topic, payload.c_str());
  }

  // MPU-6050 : envoi periodique toutes les 500 ms
  if (millis() - lastMpuSend >= MPU_INTERVAL) {
    lastMpuSend = millis();
    float accelValue = readMPUAccel();
    bool fallen = (abs(accelValue) > fall_threshold);

    JsonDocument doc;
    doc["type"]   = "gyro";
    doc["accel"]  = round(accelValue * 100) / 100.0;
    doc["fallen"] = fallen;
    doc["ts"]     = millis();
    String payload;
    serializeJson(doc, payload);
    client.publish(mqtt_topic, payload.c_str());
  }
  delay(50);
}
\end{lstlisting}

\subsection{Configuration Matérielle}

\begin{table}[H]
  \centering
  \caption{Mapping des broches ESP32}
  \label{tab:pinout}
  \renewcommand{\arraystretch}{1.4}
  \begin{tabularx}{\textwidth}{|p{2.5cm}|p{2cm}|X|}
    \hline
    \rowcolor{isimablue!20}
    \textbf{Composant} & \textbf{Broche} & \textbf{Rôle} \\
    \hline
    MPU-6050 SDA    & GPIO 21 & Ligne de données I2C \\
    \hline
    MPU-6050 SCL    & GPIO 22 & Ligne d'horloge I2C \\
    \hline
    Tamper siège    & GPIO 33 & Entrée numérique (INPUT\_PULLUP) \\
    \hline
    Tamper capot av & GPIO 25 & Entrée numérique (INPUT\_PULLUP) \\
    \hline
    Tamper capot inf& GPIO 26 & Entrée numérique (INPUT\_PULLUP) \\
    \hline
    LED indicateur  & GPIO 13 & Sortie numérique --- état tamper \\
    \hline
  \end{tabularx}
\end{table}

\begin{figure}[H]
  \centering
  \includegraphics[width=0.75\textwidth]{image_rapport/esp32_wiring.jpg}
  \caption{Câblage du prototype ESP32}
  \label{fig:wiring}
\end{figure}

%-------------------------------------------------------------------
\section{Réalisation de l'Application Mobile}
%-------------------------------------------------------------------

\subsection{Structure du Projet React Native}

\begin{lstlisting}[language=bash, caption={Structure du projet React Native / Expo}]
scooter-monitor/
|-- App.js                    # Point d'entree, gestion session
|-- src/
|   |-- screens/
|   |   |-- AuthScreen.jsx    # Ecran de connexion (invitation uniquement)
|   |   |-- HomeScreen.jsx    # Dashboard principal
|   |   |-- AdminScreen.jsx   # Gestion des comptes (admin)
|   |   |-- BatteryStockScreen.jsx  # Surveillance batteries
|   |   |-- TpmsStockScreen.jsx     # Surveillance pneus TPMS
|   |   |-- NotificationsScreen.jsx # Alertes et notifications
|   |   |-- DashboardScreen.jsx     # Vue detaillee scooter
|   |   `-- PendingScreen.jsx       # Acces non autorise
|   |-- components/
|   |   `-- Sidebar.jsx       # Menu lateral navigation
|   |-- lib/
|   |   |-- supabaseClient.js # Client Supabase configure
|   |   `-- notifications.js  # Gestion notifications
|   `-- constants.js          # Couleurs, polices, helpers
`-- supabase/
    `-- migration_*.sql       # Migrations base de donnees
\end{lstlisting}

\subsection{Intégration Supabase Realtime}

\begin{lstlisting}[language=JavaScript, caption={Abonnement temps réel Supabase (React Native)}]
useEffect(() => {
  const channel = supabase
    .channel('telemetry-changes')
    .on('postgres_changes',
      { event: 'INSERT', schema: 'public',
        table: 'telemetry',
        filter: `scooter_id=eq.${scooterId}` },
      (payload) => setTelemetry(payload.new)
    )
    .subscribe();

  return () => supabase.removeChannel(channel);
}, [scooterId]);
\end{lstlisting}

\subsection{Interfaces de l'Application Mobile}

\subsubsection{Écran d'Authentification}

Nous avons implémenté un système d'accès sur invitation uniquement~: seul l'administrateur
(compte \texttt{info@evemobility.tn}) peut créer des comptes utilisateurs, directement
depuis l'écran d'administration de l'application. L'écran de connexion n'expose aucun
lien d'inscription afin d'éviter toute création de compte non contrôlée.

\begin{figure}[H]
  \centering
  \includegraphics[width=0.35\textwidth]{image_rapport/screen_login.png}
  \caption{Écran d'authentification de l'application mobile}
  \label{fig:screen_login}
\end{figure}

\subsubsection{Dashboard Principal}

Le tableau de bord est l'écran central de l'application. Nous l'avons conçu avec EVE
Mobility pour qu'un opérateur puisse évaluer l'état d'un scooter en moins de 5 secondes
de lecture~: niveaux SOC des deux batteries avec code couleur, état de l'alarme,
indicateur de pression des pneus, et compteur d'alertes actives non archivées. Toutes
les données sont mises à jour automatiquement par Supabase Realtime sans aucun
rechargement manuel.

\begin{figure}[H]
  \centering
  \includegraphics[width=0.35\textwidth]{image_rapport/screen_dashboard.png}
  \caption{Dashboard principal --- Vue temps réel du scooter}
  \label{fig:screen_dashboard}
\end{figure}

\subsubsection{Écran de Surveillance des Batteries}

Cet écran a nécessité le plus de travail de décodage BLE~: le protocole GATT de chaque
module BMS lithium étant propriétaire, nous avons dû capturer les trames avec un
analyseur Bluetooth et rétro-ingénierié l'encodage des valeurs (tension, courant,
température, SOC, SOH). Le résultat affiché~: les valeurs numériques de chaque batterie
avec un indicateur de présence (connectée / proche / perdue).

\begin{figure}[H]
  \centering
  \includegraphics[width=0.35\textwidth]{image_rapport/screen_batteries.png}
  \caption{Écran de surveillance des batteries}
  \label{fig:screen_batteries}
\end{figure}

\subsubsection{Écran de Gestion des Alertes}

L'écran des alertes liste les événements classés par sévérité (critique, avertissement,
information). L'utilisateur peut archiver les alertes traitées.

\begin{figure}[H]
  \centering
  \includegraphics[width=0.35\textwidth]{image_rapport/screen_alerts.png}
  \caption{Écran de gestion des alertes}
  \label{fig:screen_alerts}
\end{figure}

\subsubsection{Écran de Configuration des Seuils}

\begin{figure}[H]
  \centering
  \includegraphics[width=0.35\textwidth]{image_rapport/screen_settings.png}
  \caption{Écran de configuration des seuils d'alerte}
  \label{fig:screen_settings}
\end{figure}

%-------------------------------------------------------------------
\section{Tests et Validation}
%-------------------------------------------------------------------

\subsection{Tests du Firmware ESP32}

\begin{table}[H]
  \centering
  \caption{Résultats des tests MQTT}
  \label{tab:test_mqtt}
  \renewcommand{\arraystretch}{1.4}
  \small
  \begin{tabularx}{\textwidth}{|X|c|c|}
    \hline
    \rowcolor{isimablue!20}
    \textbf{Cas de test} & \textbf{Attendu} & \textbf{Obtenu} \\
    \hline
    Connexion WiFi (SSID valide) & Connexion établie & \textcolor{green!60!black}{✓ Succès} \\
    \hline
    Connexion MQTT HiveMQ (port 8883) & rc = 0 & \textcolor{green!60!black}{✓ Succès} \\
    \hline
    Publication message JSON (gyro) & Message reçu & \textcolor{green!60!black}{✓ Succès} \\
    \hline
    Publication message JSON (contact) & Message reçu & \textcolor{green!60!black}{✓ Succès} \\
    \hline
    Reconnexion automatique après coupure & Reprise en $<$10~s & \textcolor{green!60!black}{✓ Succès} \\
    \hline
    Connexion WiFi (SSID invalide) & Échec gracieux & \textcolor{green!60!black}{✓ Succès} \\
    \hline
  \end{tabularx}
\end{table}

\begin{table}[H]
  \centering
  \caption{Résultats des tests MPU-6050}
  \label{tab:test_mpu}
  \renewcommand{\arraystretch}{1.4}
  \small
  \begin{tabularx}{\textwidth}{|p{3.5cm}|p{3cm}|p{3cm}|c|}
    \hline
    \rowcolor{isimablue!20}
    \textbf{Position} & \textbf{Attendu} & \textbf{Obtenu} & \textbf{Résultat} \\
    \hline
    À plat (horizontal) & accel\_z $\approx$ 90° & 88,4° & \textcolor{green!60!black}{✓} \\
    \hline
    Incliné 45° (gauche) & accel\_x $\approx$ 45° & 43,7° & \textcolor{green!60!black}{✓} \\
    \hline
    Retourné (180°) & accel\_z $\approx$ -90° & -88,9° & \textcolor{green!60!black}{✓} \\
    \hline
    Debout vertical & accel\_x, z $\approx$ 0° & 1,2° ; 0,8° & \textcolor{green!60!black}{✓} \\
    \hline
  \end{tabularx}
\end{table}

\begin{table}[H]
  \centering
  \caption{Résultats des tests capteur de contact (tamper)}
  \label{tab:test_contact}
  \renewcommand{\arraystretch}{1.4}
  \small
  \begin{tabularx}{\textwidth}{|X|c|c|}
    \hline
    \rowcolor{isimablue!20}
    \textbf{Cas de test} & \textbf{Attendu} & \textbf{Obtenu} \\
    \hline
    Contact fermé (capot fermé) & value = 1, LED allumée & \textcolor{green!60!black}{✓ Succès} \\
    \hline
    Contact ouvert (capot ouvert) & value = 0, LED éteinte & \textcolor{green!60!black}{✓ Succès} \\
    \hline
    Publication MQTT lors du changement & Message immédiat & \textcolor{green!60!black}{✓ Succès} \\
    \hline
    Pas de publication si état inchangé & Aucun message & \textcolor{green!60!black}{✓ Succès} \\
    \hline
  \end{tabularx}
\end{table}

\subsection{Tests d'Intégration Bout-en-Bout}

\begin{table}[H]
  \centering
  \caption{Résultats des tests d'intégration bout-en-bout}
  \label{tab:test_integration}
  \renewcommand{\arraystretch}{1.5}
  \small
  \begin{tabularx}{\textwidth}{|X|p{3cm}|c|}
    \hline
    \rowcolor{isimablue!20}
    \textbf{Scénario} & \textbf{Attendu} & \textbf{Statut} \\
    \hline
    Données gyroscope ESP32 $\to$ Supabase $\to$ App & Affichage en $<$2~s & \textcolor{green!60!black}{✓ Validé} \\
    \hline
    Détection ouverture capot $\to$ alerte App & Notification en $<$1~s & \textcolor{green!60!black}{✓ Validé} \\
    \hline
    Commande alarme App $\to$ ESP32 & Exécution en $<$3~s & \textcolor{green!60!black}{✓ Validé} \\
    \hline
    Reconnexion après perte réseau & Reprise automatique & \textcolor{green!60!black}{✓ Validé} \\
    \hline
    Authentification utilisateur & Accès dashboard sécurisé & \textcolor{green!60!black}{✓ Validé} \\
    \hline
    Archivage d'une alerte & Disparaît de la liste active & \textcolor{green!60!black}{✓ Validé} \\
    \hline
  \end{tabularx}
\end{table}

\subsection{Tests de Performance}

\begin{table}[H]
  \centering
  \caption{Mesures de performance du système}
  \label{tab:test_perf}
  \renewcommand{\arraystretch}{1.4}
  \small
  \begin{tabularx}{\textwidth}{|X|p{3cm}|p{3.2cm}|}
    \hline
    \rowcolor{isimablue!20}
    \textbf{Indicateur} & \textbf{Objectif} & \textbf{Résultat} \\
    \hline
    Latence MQTT (ESP32 $\to$ Supabase) & $<$500~ms & 120~ms en moyenne \\
    \hline
    Latence affichage App & $<$2~s & 800~ms en moyenne \\
    \hline
    Disponibilité Supabase & $>$99,5\,\% & 99,9\,\% \\
    \hline
    Taux de perte de messages MQTT & $<$1\,\% & 0,2\,\% \\
    \hline
    Fréquence publication MPU-6050 & 2~msg/s & 1,98~msg/s \\
    \hline
    Mémoire heap ESP32 libre & $>$30~kB & 42~kB en régime permanent \\
    \hline
  \end{tabularx}
\end{table}

%-------------------------------------------------------------------
\section*{Conclusion}
\addcontentsline{toc}{section}{Conclusion}
%-------------------------------------------------------------------

\begin{chapterconclusion}
Ce chapitre a présenté la réalisation complète du système. Le firmware ESP32 assure
la collecte et la transmission fiable des données capteurs via MQTT over TLS, avec
une latence moyenne de 120~ms. L'application mobile React Native, intégrée à Supabase
Realtime, offre une supervision intuitive et réactive. Les tests d'intégration ont
confirmé la robustesse du système~: disponibilité de 99,9\,\%, taux de perte MQTT
de 0,2\,\% et détection fiable de tous les événements critiques testés.
\end{chapterconclusion}

%=====================================================================
% CONCLUSION GÉNÉRALE
%=====================================================================
\clearpage
\chapter*{Conclusion Générale}
\addcontentsline{toc}{chapter}{Conclusion Générale}
\markboth{Conclusion Générale}{Conclusion Générale}
\thispagestyle{fancy}

Au terme de ce projet de fin d'études mené chez EVE Mobility, nous livrons un système
de télésurveillance fonctionnel, testé sur le matériel réel de la société, et prêt à
être déployé en production sur la flotte.

Le bilan technique est le suivant~: le firmware ESP32 transmet correctement les données
de trois catégories de capteurs (MPU-6050, détecteurs tamper, TPMS) vers Supabase via
HiveMQ Cloud avec une latence mesurée de 120~ms en moyenne. L'application React Native
affiche ces données en temps réel grâce aux abonnements Supabase Realtime, et permet
à l'administrateur de gérer les comptes utilisateurs par invitation. La base de données
PostgreSQL, organisée en six tables (\texttt{scooters}, \texttt{telemetry},
\texttt{batteries}, \texttt{tpms\_sensors}, \texttt{profiles}, \texttt{invited\_emails}),
est sécurisée par des politiques Row Level Security qui garantissent la séparation
stricte des données entre opérateurs.

Sur le plan humain, ce stage nous a confrontés à des défis réels que les exercices
académiques ne reproduisent pas~: rétro-ingénierie du protocole BLE propriétaire des
modules BMS, gestion des reconnexions MQTT après coupure Wi-Fi, et conception d'une
interface suffisamment intuitive pour être utilisée sans formation par les opérateurs
d'EVE Mobility. Ces difficultés ont été nos meilleurs enseignements.

Plusieurs axes d'amélioration ont déjà été identifiés avec l'équipe d'EVE Mobility~:

\begin{itemize}
  \item \textbf{Enrichissement des données batteries}~: le champ \texttt{soc} pourrait
        être complété par la tension, le courant et la température, nécessitant un
        décodage plus complet du protocole BMS.
  \item \textbf{Mode deep sleep ESP32}~: mise en veille périodique du microcontrôleur
        pour réduire la consommation sur les scooters en stationnement prolongé.
  \item \textbf{Filtre de Kalman sur la détection de chute}~: réduire les faux
        positifs générés sur les routes mal entretenues tout en maintenant la
        sensibilité aux chutes réelles.
  \item \textbf{Notifications push}~: intégration d'Expo Notifications pour alerter
        l'opérateur même lorsque l'application est en arrière-plan.
  \item \textbf{Tableau de bord multi-scooters}~: vue synthétique de toute la flotte
        sur un seul écran, avec filtres par état (\texttt{online}, \texttt{fallen},
        \texttt{alarm}).
\end{itemize}

Ce projet nous a montré concrètement comment une architecture IoT moderne --- du capteur
à l'écran du téléphone --- peut transformer la gestion opérationnelle d'une flotte de
véhicules électriques. Nous sommes fiers de remettre à EVE Mobility un outil qu'ils
pourront faire évoluer bien au-delà du cadre de ce stage.

%=====================================================================
% BIBLIOGRAPHIE
%=====================================================================
\clearpage
\begin{thebibliography}{99}
  \addcontentsline{toc}{chapter}{Bibliographie}
  \markboth{Bibliographie}{Bibliographie}

  \bibitem{esp32}
  Espressif Systems,
  \textit{ESP32 Technical Reference Manual},
  version 5.x, 2023.
  Disponible sur~: \url{https://www.espressif.com/sites/default/files/documentation/esp32_technical_reference_manual_en.pdf}

  \bibitem{supabase}
  Supabase Inc.,
  \textit{Supabase Documentation --- Database, Auth, Realtime, Storage},
  2024.
  Disponible sur~: \url{https://supabase.com/docs}

  \bibitem{reactnative}
  Meta Platforms,
  \textit{React Native Documentation},
  version 0.73, 2024.
  Disponible sur~: \url{https://reactnative.dev/docs/getting-started}

  \bibitem{expo}
  Expo Dev,
  \textit{Expo Documentation --- Build, Deploy, Update},
  2024.
  Disponible sur~: \url{https://docs.expo.dev}

  \bibitem{mpu6050}
  InvenSense (TDK),
  \textit{MPU-6050 Product Specification, Revision 3.4},
  2013.
  Disponible sur~: \url{https://invensense.tdk.com/wp-content/uploads/2015/02/MPU-6000-Datasheet1.pdf}

  \bibitem{ble}
  Bluetooth SIG,
  \textit{Bluetooth Core Specification 5.3},
  2021.
  Disponible sur~: \url{https://www.bluetooth.com/specifications/specs/core-specification-5-3/}

  \bibitem{mqtt}
  OASIS Standard,
  \textit{MQTT Version 5.0 --- Message Queuing Telemetry Transport},
  2019.
  Disponible sur~: \url{https://docs.oasis-open.org/mqtt/mqtt/v5.0/mqtt-v5.0.html}

  \bibitem{hivemq}
  HiveMQ GmbH,
  \textit{HiveMQ Cloud Documentation},
  2024.
  Disponible sur~: \url{https://www.hivemq.com/docs/hivemq-cloud/}

\end{thebibliography}

\end{document}
